\documentclass[a4paper, 12pt]{article}



\usepackage[utf8]{inputenc}
\usepackage[T1]{fontenc}
\usepackage{textcomp}
\usepackage{amssymb}
\usepackage{newtxtext} \usepackage{newtxmath}
\usepackage{amsmath, amssymb}
\newtheorem{problem}{Problem}
\newtheorem{example}{Example}
\newtheorem{lemma}{Lemma}
\newtheorem{theorem}{Theorem}
\newtheorem{problem}{Problem}
\newtheorem{example}{Example} \newtheorem{definition}{Definition}
\newtheorem{lemma}{Lemma}
\newtheorem{theorem}{Theorem}
\DeclareMathAlphabet{\mathcal}{OMS}{cmsy}{m}{n}
\usepackage{parskip}

\begin{document}

    
$\S$ Il est une chose curieuse : avant, je pensais que le meilleur de moi
fleurissait quand j'étais en compagnie des autres. Seul, je me découvrais
mélancolique, morne et taciturne. Mais, alors, je comprends que précisément
lorsque je suis seul, et précisément depuis cette bataille contre ma tristesse,
ma vie dégage sa lumière la plus salutaire et secrète. 

La bataille, il faut bien le dire, est cruelle, et peut-être qu'il y a un
certain profit dans la considération de sa nature. Mais, il n'est pas clair par
où on doit commencer : peut-être, je peux commencer ainsi : Il n'y a rien de
plus étrange pour moi que ceux qui sont à la fois heureux et simples. Même quand
je suis heureux, cet état n'est qu'un équilibre fragile, comme celui d'une boule
sur le maximum local d'une courbe concave, au point où la dérivée est exactement
zéro. Les forces du monde, extérieures et intérieures, l'incitent à trembler,
bien que dans un rayon limité, violemment. Et il est impossible de nier que ce
genre de personnes existe : des gens de bien qui ne sont pas assiégés par le
doute, l'agonie, l'affrontement entre le bien et le mal. J'ai pensé parfois, en
les regardant : « Le pauvre ! Qu'est-ce que tu feras quand la vie ouvre sa
gueule insensible devant ton innocence ? » Mais c'est une question idiote :
qu'est-ce que je fais ? En outre, cette classe de personnes, il faut le dire,
est belle, même terriblement. Ce qui me perturbe est simplement que, face à
leurs visages, je découvre que mon propre visage est une toile décatie sur
laquelle la vie a gravé un dessein jamais désiré par moi. Notez que je ne dis
pas, et ne prétends jamais dire, que je ne suis pas heureux : plutôt, je veux
dire simplement que le bonheur est pour moi comme l'eau coulante et éphémère
d'un ruisseau, plutôt que l'impassible surface d'un lac. 


C'est alors ce que je veux dire : les personnes qui sont à la fois simples et
heureuses me font peur. Devant leur hédonisme ---elles ne peuvent être
qu'hédonistes--- leur courage, leur insouciance, je vois le pire en moi et je
sens comment mon rêve idiot d'être compris s'écroule. Comme je peux être heureux
avec une âme déchirée ; non pas misérable, mais déchirée ! Mais face à un de ces
cerfs immaculés et vierges, quel gouffre, quel reflet infernel, quelles ténèbres
!


Ce que je dis n'est pas le point central de mon discourse, mais il sert à
suggérer dans quelle mesure le bonheur est pour moi une bataille sanglante.
Naturellement, ce n'est pas une chose spéciale : c'est ainsi pour la plupart des
hommes. Comment se fait-il que les hommes puissent vivre, que je puisse vivre,
dans cette bataille ? 

\begin{quote} La beauté est aussi mystérieuse que terrible. Dieu et le diable
s'y affrontent, et le champ de bataille est le cœur de l'homme \end{quote}

Dans \textit{L'idiot}, Dostoïevski explore l'idée selon laquelle la beauté peut
sauver le monde. La beauté signifie ici la forme la plus parfaite du
christianisme. Mais il est clair que la conclusion du roman n'est pas optimiste
--- c'est le moins qu'on puisse dire. Mais le salut du monde est une ambition
très grande : le salut d'un jour, la transformation d'un jour affaibli en un
jour heureux, est possible seulement à travers la beauté. La beauté est une
chose difficile à comprendre pour moi, car je ne suis pas un Alyosha, mais un
Ivan --- pour ainsi dire. Je déteste mon cerveau. Il est clair que si je suis
abattu et que j'écris, je me retrouve pareil ou pire ; mais si j'embrasse
quelqu'un que j'aime, si j'aide quelqu'un qui en a besoin, si je suis comme
l'homme gentil qu'Antonio Machado a imaginé, et que je bois du vin s'il y en a,
et sinon de l'eau fraîche, mon cœur devient sain et la lumière du monde brûle en
lui.

Ma tristesse n'est pas liée à circonstances de ma vie : au contraire, ma vie est
un privilège et rien de plus. Oui, certaines blessures du passé sont latentes
encore : un suicide, un mensonge, une famillie. Mais je ne suis pas victime de
cela. Ma souffrance est presque intellectuel. J'allais en bus hier : en
regardant les rues sordides et sales, je pensais que peut-être mon souffrir est
la voix de Dieu m'appelant, je me rappelais l'origine de mon nom et l'homme qui
a lutté avec Dieu, car je flairais une chose étrange et spirituelle crepiter en
moi. Soudainement, un vieille homme humble, qu'était assis près de moi, me
regarde et dit : «Dis-moi : quel est le secret pour ne pas vieillir !». Alors,
il me raconte l'histoire triste de sa pauvreté, très habituelle dans ses traits
généraux, et, en descendant du bus, termine en disant : «Toi, étudie pour faire
la révolution !». Silencieux, je pensai : «Bon monsieur, j'abandonne
l'Argentine, je vais en Suisse, et ce qui est pire, moi qu'ai tout, je ne suis
pas heureux !».  

Quelques heures plus tard je suis arrivé à ma destination : la maison abandonnée
d'un oncle éloigné, un homme schizophrène qui est décédé, et que je garde
quelques jours par semaine. La maison est pleine de grilles, de fils barbelés,
de morceaux de verre, et de cadenas, car cet oncle avait une obsession paranoïde
pour la sécurité. À cette occasion, j'étais là-bas pour accueillir un homme qui
avait la tâche de réparer certains problèmes dans la canalisation de la maison.
Cet homme, qui avait soixante-dix ans, plus ou moins, qui était douce, gentil,
et enclin à faire des blagues inappropriées, montait péniblement l'échelle
jusqu'au toit et se plaignait du bruit d'une tondeuse à gazon en criant : « j'ai
des acouphènes !» Ses yeux étaient pleins de lumière. La plupart de sa retraite
était consacrée au traitement de son chien, qui était tombé malade du jour au
lendemain, et, en outre, son gendre lui demandait de l'argent pour
je-ne-sais-quoi. Toutefois, son sourire était affectueux et il jurait toutes les
cinq minutes qu'il ne perdait pas la foi en Dieu.

Alors, ma journée touchait à sa fin. Cet écrit est chaotique, car mes pensées le
sont. J'ai commencé en parlant de ma solitude, d'une bataille de laquelle une
fleur fleurissait, et lentement je suis arrivé à la considération de certains
personnages, à des choses très concrètes, très réelles. C'est précisément mon
problème : mais voilà alors que, quand mes obligations se sont terminées, et que
j'étais seul dans la nuit comme un chien, en rentrant chez moi, j'ai vu un
cartonnier --- un «ciruja» --- cherchant dans les ordures. Je l'ai nourri et je
lui ai demandé son histoire. Cet homme était jeune, et ses grands yeux étaient
lumineux ; c'était un homme du Christ et il me parlait de la lutte du bien et du
mal --- il a utilisé ces termes ---, du fait que peu importe à quel point on
essaie, le bien est une étoile très difficile à atteindre, et de la manière dont
la police ne lui permet pas de travailler en paix. La police, m'a-t-il dit,
après lui avoir demandé son nom, l'appelle de toute façon « negro » : «ils ne
m'appellent pas Rodrigo, ou monsieur Hernández, mais simplement noir». Il m'a
dit qu'il était allé en prison, qu'il n'allait pas à l'église, car «avec toutes
les choses qu'il a vues» sa spiritualité est intérieure. Enfin, nous avons parlé
comme ce que nous étions : deux hommes inconnus qui ne se retrouveront plus. Et
il est clair qu'aucune des pensées tordues que j'ai eues durant la journée ne
vaut une seule seconde passée avec cet homme, ni sa sasiété, ni le sentiment non
médiatisé par la pensée qui surgit quand on est en communion avec les autres
hommes.



























\end{document}



